\section{Selective Boundary Refinement}
\label{sec:method}
\subsection{Problem and incumbent}
Given audio $x$ of duration $D$ and query $q$, let the incumbent prediction be $S_0=\{[s_i,e_i]\}_{i=1}^{K}$. Intervals are ordered, valid, and non-overlapping. For a ground-truth interval set $G$, temporal set-IoU is
\begin{equation}
\setiou(S,G)=\frac{|S\cap G|}{|S\cup G|},
\label{eq:iou}
\end{equation}
where $|\cdot|$ denotes temporal measure of the interval union. The refinement stage changes boundaries, not event identity or count. An empty incumbent remains empty.

The full system reuses an existing strong predictor rather than generating a new answer from scratch. On SpotSound-Bench, its incumbent is the locked E004 NOVA configuration: relation-aware adaptation and setwise preference optimization (SetPO) provide candidate predictions, and keep/drop audio verification supports conservative candidate selection. On the other reported grounding sets, the incumbent is the official SpotSound-A prediction. This distinction is retained throughout evaluation.

\subsection{Legal edits and relative utility}
For event $i$, candidate endpoints lie on a fine acoustic grid within
\begin{equation}
r_i=\min(r_{\max},\rho(e_i-s_i)).
\label{eq:radius}
\end{equation}
The original endpoints are also included exactly. Each option $a_i$ must remain within $[0,D]$, span at least one grid step, and respect neighboring incumbent intervals. Invalid candidate construction causes abstention. Keeping all intervals unchanged is always an available action.

Frozen SetPO/SpanTool evidence provides occupancy, onset, offset, and fine boundary channels. The edit feature vector $\phi_i(a_i)$ contains changes in acoustic boundary signatures, normalized endpoint displacements, and geometry--context interactions involving duration, coverage, and neighboring gaps. It is constructed so that $\phi_i(\keep)=0$ exactly.

Training labels measure the effect of editing one interval while leaving all others unchanged:
\begin{equation}
y_{n,i,a}=
\setiou(S_n^{i\leftarrow a},G_n)
-\setiou(S_n,G_n).
\label{eq:target}
\end{equation}
A ridge model predicts this relative utility. With feature-wise root-mean-square scaling $C$ and normalized weights $\omega_j$, it solves
\begin{equation}
\hat w=\arg\min_w
\sum_j\omega_j\big(w^\top C^{-1}\phi_j-y_j\big)^2
+\alpha\|w\|_2^2.
\label{eq:ridge}
\end{equation}
Weights balance data sources, their audio groups, the queries in each group, and the options for each query. No intercept is fitted; the identity utility is exactly zero. The acoustic backbone is frozen during this stage, but utility fitting is supervised: the complete approach is not training-free.

\subsection{Structured selection and abstention}
Let $\mathcal A_i$ be event $i$'s legal options. We select one option per event to maximize the additive learned surrogate,
\begin{equation}
a^\star=\arg\max_{a\in\mathcal F}
\sum_{i=1}^{K}\hat w^\top C^{-1}\phi_i(a_i),
\label{eq:select}
\end{equation}
where $\mathcal F$ also prohibits overlap between selected intervals. The dynamic-programming state for option $a$ at layer $i$ is
\begin{equation}
V_i(a)=u_i(a)+
\max_{\substack{b\in\mathcal A_{i-1}\\e(b)<s(a)}}V_{i-1}(b).
\label{eq:dp}
\end{equation}
Options are sorted by endpoint and compatible-prefix maxima are reused. This exactly optimizes the \emph{additive utility surrogate}, not the true multi-interval set-IoU, whose changes need not be additive.

To assess stability, $B$ utility models are fitted by resampling audio groups within each source. For the selected candidate, define
\begin{equation}
L(a^\star)=Q_{0.05}\!\left(
\left\{\sum_i u_i^{(b)}(a_i^\star)\right\}_{b=1}^{B}
\right).
\label{eq:quantile}
\end{equation}
We apply the candidate only if its reference-model utility is positive and $L(a^\star)>m$; otherwise we retain $S_0$. Here $Q_{0.05}$ is an ensemble quantile, not a certified lower bound on realized test improvement. Ground-truth intervals are used for utility training and offline evaluation, never for selecting an edit at inference.
